
\section{定理 \ref{thm:wl2_to_cc} 的自足证明思路}

给定边着色有向图 $G=(V,E)$。$\WLTWO$ 在 $V\times V$ 上迭代赋色：初始颜色由有限信息决定（例如 $(u=v)$ 与边颜色/方向等），每轮更新把旧颜色与如下“证据多重集”组合哈希为新颜色：
$$
M(u,v):=\multiset{(\chi(u,w),\chi(w,v)):\ w\in V},
$$
其中 $\chi$ 表示上一轮在 $V\times V$ 上的颜色函数。

当 $\WLTWO$ 稳定时，若两对 $(u,v)$ 与 $(u',v')$ 同色，则它们对任意颜色对 $(r,s)$ 的中介计数
$$
N_{rs}(u,v):=\abs{\{w\in V:\ \chi(u,w)=r,\ \chi(w,v)=s\}}
$$
必相等；否则证据多重集 $M(u,v)$ 与 $M(u',v')$ 将在下一轮不同而被分色。于是对每个颜色三元组 $(r,s,t)$，可定义交数
$$
p_{rs}^t:=N_{rs}(u,v)\quad \text{其中 } \chi(u,v)=t,
$$
并由稳定性保证其与 $(u,v)$ 的具体选择无关，故交数良定义。

对角线分解与转置闭合来自初始颜色与更新规则在 $(u,v)\leftrightarrow (v,u)$ 以及 $u=v$ 情形下的显式对称性。综上，稳定后的 $V\times V$ 颜色划分满足相干配置的三条公理，从而得到定理 \ref{thm:wl2_to_cc}。

